\begin{theorem}[A strict gate from the odd harmonic CM image]
\label{thm:odd-harmonic-cm-gate}
Put
\[
 P_D=D\cap\bigcap_p\ker\Phi_p,\qquad
 T_D=P_D\cap\ker d_{(3)},
\]
where \(d_{(3)}\) is the cubic radical--CM defect.
Let \(W_{\rm odd}\) be the projection of the harmonic CM
homomorphism \(W:D\to N_4(4)\) onto the maximal pro-odd factor,
and put
\[
 U_D=T_D\cap\ker W_{\rm odd}.
\]
Then \(U_D\) is closed and normal in \(G\), and
\begin{equation}
 B_{\max2}\cong
 \left(U_D/\overline{[G,U_D](U_D\cap H)}\right)_{\max2}.  \label{eq:odd-cm-1}
\end{equation}
The kernel of the original restriction map is
\begin{equation}
 \ker\bigl(X_2\to H^1(U_D,\mathbf F_2)\bigr)
  =\bigoplus_{q\ {\rm odd}}\mathbf F_2\psi_q
                           \oplus\mathbf F_2\eta_3.     \label{eq:odd-cm-2}
\end{equation}
Moreover
\begin{equation}
 [T_D:U_D]=\infty,\qquad
 [T_D\cap H:U_D\cap H]=\infty.                          \label{eq:odd-cm-3}
\end{equation}
\end{theorem}

\begin{proof}
The pro-odd projection is characteristic, and the conjugation
action on \(W\) is identity or inversion. Thus \(U_D\) is
closed and \(G\)-normal. Its quotient in \(T_D\) is the
pro-odd group \(W_{\rm odd}(T_D)\).

A \(G\)-invariant character of \(U_D\) with values in a finite
abelian two-group extends uniquely to a \(G\)-invariant
character of \(T_D\).
Indeed its pushout is a central extension of a pro-odd group
by a finite two-group. In a finite quotient, if \(c\) is the
central two-cocycle and the odd quotient has order \(m\), the
cochain
\[
 b(q)=m^{-1}\sum_t c(q,t)
\]
satisfies \(c(q,t)=b(q)+b(t)-b(qt)\).
The resulting complement is unique, since an odd-order group
has no nonzero character in a two-group. Passing to the
inverse limit gives the unique extension.
Uniqueness makes it \(G\)-invariant.
If the original character kills \(U_D\cap H\), the extension
kills \(T_D\cap H\), because
\((T_D\cap H)/(U_D\cap H)\) is pro-odd.
Thus the relative maximal pro-two abelian quotients on
\(T_D\) and \(U_D\) agree.
The established reduction on \(T_D\) gives \eqref{eq:odd-cm-1}.

Likewise a character of \(T_D\) with values in \(\mathbf F_2\)
and vanishing on \(U_D\) factors through a pro-odd quotient
and is zero. The original restriction kernels on \(T_D\)
and \(U_D\) therefore agree, giving \eqref{eq:odd-cm-2}.
Neither argument assumes \(G=T_DH\).

We construct an arithmetic witness for \eqref{eq:odd-cm-3}.
The arithmetic CM image theorem and the harmonic comparison
give
\[
 \rho_4(M\cap H)=N_4(4)
\]
for the Tate representation of \(E_4:y^2=x^3-x\), and identify
its odd projection on \(D\cap H\) with \(W_{\rm odd}\).
The primary reciprocity input is
\cite[Theorem~12.27, p.~124]{MilneMF110}.
At the split prime five,
\[
 \mathbf Z[i]\otimes\mathbf Z_5
       \cong\mathbf Z_5\times\mathbf Z_5,
\]
and the norm-one torus contains \(v=(6,6^{-1})\), of
infinite pro-five order.
Place this value at five and the identity at all other
components of \(N_4(4)\). Its two-component satisfies the
level-four condition.
Choose \(m\in M\cap H\) with \(\rho_4(m)=v\), and put
\[
 a=m^{\varepsilon_5},
\]
where \(\varepsilon_5\in\widehat{\mathbf Z}\) retains the
five-component and kills all others.
Then \(a\) is pro-five, has full multiplier one, and
\(\rho_4(a)=v\).

Choose \(g\in H\) with full multiplier
\[
 \lambda_2(g)=-1,\qquad
 \lambda_p(g)=1\quad(p\text{ odd}),
\]
and set \(h=g^{\varepsilon_2}\).
Then \(h\) is pro-two and has the same multiplier tuple.
For every \(p\), the action image of \(\ker\lambda_p\) is
pro-\(p\), by its upper-unitriangular action on finite
mod-\(p\) Magnus truncations.
It follows that \(\Phi_p(h)=1\) for every odd \(p\):
its procyclic image is both pro-two and pro-\(p\).
Similarly \(\Phi_p(a)=1\) for \(p\ne5\).

Consequently the arithmetic commutator
\[
                         w=[h,a]
\]
has trivial \(\Phi_p\)-image for every prime \(p\).
Its full multiplier is one.
Every rational Kummer value of \(a\) is supported at five,
where \(\lambda_5(h)=1\), so the conjugation law gives
\[
 \rho_q^I(w)
       =(\lambda(h)-1)\rho_q^I(a)=0
       \qquad(q\text{ a rational prime}).
\]
Thus \(w\in P_D\cap H\).
The cubic defect vanishes on \(D\cap H\), hence
\(w\in T_D\cap H\).

Finally \(h\) acts nontrivially on \(\mathbf Q(i)\), because
its dyadic multiplier is \(-1\) modulo four.
It conjugates the CM coefficient ring by complex conjugation
and therefore inverts its norm-one torus.
Compute in the Tate representation, so that no value
of \(W\) outside its stated domain is needed:
\[
                    W_{\rm odd}(w)=v^{-1}v^{-1}=v^{-2}.
\]
This element still has infinite pro-five order.
The procyclic closure of \(w\) lies in \(T_D\cap H\) and has
infinite image in the quotient by \(U_D\).
Both assertions in \eqref{eq:odd-cm-3} follow.
\end{proof}

\begin{remark}
The strictness witness simultaneously kills every primary
fundamental-group action and every rational Kummer coordinate.
It requires no joint radical--CM image assertion.
The theorem does not claim \(G=U_DH\) or compute the remaining
relative abelian quotient.
\end{remark}

\begin{corollary}\label{cor:full-odd-harmonic-arithmetic-image}
Let \(A_{4,\rm odd}=N_4(4)_{\rm odd}\). Then
\[
 W_{\rm odd}(T_D\cap H)=W_{\rm odd}(T_D)=A_{4,\rm odd},
 \qquad T_D/U_D\cong A_{4,\rm odd}.
\]
Moreover \(T_D=U_D(T_D\cap H)\).
This factorization is internal to \(T_D\).
\end{corollary}
\begin{proof}
For any \(v\in A_{4,\rm odd}\), the arithmetic CM image gives
\(m\in M\cap H\) with odd CM value \(v\).
Put \(a=m^{1-\varepsilon_2}\). Its procyclic closure is pro-odd,
its full multiplier is one, its odd CM value is \(v\), and
its rational Kummer values have zero dyadic component.
In particular \(\Phi_2(a)=1\).
Use the same arithmetic pro-two element \(h\) as in the theorem.
It has trivial \(\Phi_p\)-image for every odd \(p\).
Thus \(w=[h,a]\) has trivial image at every prime.
The Kummer identity
\[
 \rho_q^I(w)=(\lambda(h)-1)\rho_q^I(a)=0
\]
holds componentwise: the dyadic coordinate of \(a\) is zero,
and the odd multiplier coordinates of \(h\) are one.
Therefore \(w\in P_D\cap H\), and arithmetic vanishing of
\(d_{(3)}\) puts it in \(T_D\cap H\).
The same CM conjugation calculation gives
\(W_{\rm odd}(w)=v^{-2}\).
Raising to \(-2\) is an automorphism of the profinite abelian
pro-odd group \(A_{4,\rm odd}\), so this proves the arithmetic
surjectivity. Compactness gives the asserted quotient.
For \(t\in T_D\), choose \(j\in T_D\cap H\) with the same
odd CM value. Then \(tj^{-1}\in U_D\), proving the factorization.
\end{proof}
